\chapter{Solutions to Selected Exercises}\label{ch:solutions}
\section{Chapter~\ref{ch:modform}}%ack: r. bradshaw wrote all these!
\begin{enumerate}
\item\exref{ch:modform}{ex:upperhalfpres}.
Suppose $\gamma = \abcd{a}{nb}{c}{d} \in \GL_2(\R)$ is a matrix
with positive determinant.
Then $\gamma$ is a linear fractional transformation that fixes the
real line, so it must either fix or swap the upper and
lower half planes. Now
$$\gamma(i) = \frac{ai+b}{ci+d} = \frac{ac+bd + (ad-bc)i}{d^2+c^2},$$
so since $\det \gamma = ad-bc > 0$, the imaginary part of $\gamma(i)$
is positive; hence $\gamma$ sends the upper half plane to itself.

\item\exref{ch:modform}{ex:meromorphic}.
Avoiding poles, the quotient rule for differentiation goes
  through exactly as in the real case, so any rational
  function $f(z) = p(z)/q(z)$ ($p,q \in \C[z]$) is holomorphic on
  $\C-\{\alpha : q(\alpha) = 0 \}$. By the fundamental theorem of
  algebra, this set of poles is finite, and hence it is discrete.
  Write $q (z) = a_n(z-\alpha_1)^{r_1}\cdots(z-\alpha_k)^{r_k}$ for
  each $\alpha_i$ and let $q_i(z)=q(z)/(z-\alpha_i)^{r_i}$ which is a
  polynomial nonzero at $\alpha_i$. Thus for each $i$ we have
  $(z-\alpha_i)^{r_i}f(z) = p(z)/q'(z)$ is holomorphic at $\alpha_i$
  and hence $f(z)$ is meromorphic on $\C$.

\item \exref{ch:modform}{ex:wmfprod}.
\begin{enumerate}
\item The product $fg$ of two meromorphic functions on the upper half
  plane is itself meromorphic. Also, for all $\gamma \in \SL_2(\Z)$
  we have
\begin{align*}
(fg)^{[\gamma]_{k+j}} 
&= \frac{1}{(cz+d)^{k+j}}((fg) \circ \gamma)\\
&= \frac{1}{(cz+d)^k}(f \circ \gamma) \frac{1}{(cz+d)^j}(g \circ \gamma)
= fg,
\end{align*} so $fg$ is weakly modular. 

\item If $f$ is meromorphic on the upper half plane, then so is $1/f$.
 Now
$$
  \qquad\qquad\frac{1}{f} = \frac{1}{(cz+d)^{-k} f \circ \gamma} = (cz+d)^k ((1/f) \circ \gamma) = \frac{1}{f}^{[\gamma]_{-k}},
$$
so $1/f$ is a weakly modular form of weight $-k$. 

\item Let $f$ and $g$ be modular functions. Then, as above, $fg$ is a
  weakly modular function.  Let $\sum_{n=m}^\infty a_nq^n$ and
  $\sum_{n=m'}^\infty b_nq^n$ be their $q$-expansions around any
  $\alpha \in \P^1(\Q)$; then their formal product is the
  $q$-expansion of $fg$. But the formal product of two Laurent series
  about the same point is itself a Laurent series with convergence in
  the intersection of the convergent domains of the original series,
  so $fg$ has a meromorphic $q$-expansion at each $\alpha \in
  \P^1(\Q)$ and hence at each cusp.

\item We are in exactly the same case as in part (c), but because $f$
  and $g$ are modular functions, $m, m'\ge 0$ and hence the function
  is holomorphic at each of its cusps.

\end{enumerate}

\item \exref{ch:modform}{ex:nomodformodd}.
Let $f$ be a weakly modular function of odd weight
$k$. Since
$\gamma 
   = \abcd{-1}{0}{0}{-1} \in \SL_2(\Z)$, we have $f(z) =
(-1)^{-k}f(\gamma(z)) = -f(z)$ so $f = 0$.

\item  \exref{ch:modform}{ex:allsame}.
  Because $\SL_2(\Z/1\Z)$ is the trivial group, $\Gamma(1) =
  \ker(\SL_2(\Z) \rightarrow \SL_2(\Z/1\Z))$ must be all of
  $\SL_2(\Z)$.  As $\SL_2(\Z) = \Gamma(1) \subset \Gamma_1(1)
  \subset \Gamma_0(1) \subset \SL_2(\Z)$, we must have $\Gamma(1) =
  \Gamma_1(1) = \Gamma_0(1) = \SL_2(\Z)$.

\item \exref{ch:modform}{ex:conggamma1}.
\begin{enumerate}
\item The group $\Gamma_1(N)$ is the inverse image of
the subgroup of $\SL_2(\Z/N\Z)$ generated by $\abcd{1}{1}{0}{1}$, and
the inverse image of a group (under a group homomorphism)
is a group.
\item The group contains the kernel of the homomorphism
$\SL_2(\Z) \to \SL_2(\Z/N\Z)$, and 
that kernel has finite index since the quotient is
contained in $\SL_2(\Z/N\Z)$, which is finite. 
\item Same argument as previous part. 
\item The level is at most $N$ since both groups contain $\Gamma(N)$.
  It can be no greater than $N$ since $\abcd{1}{N}{0}{1}$ is in both
  groups.
\end{enumerate}

\item \exref{ch:modform}{ex:grpact}.
See \cite[Lemma~1.2.2]{diamond-shurman}.

\item \exref{ch:modform}{ex:sl2ztrans}.
   Let $\alpha = p/q \in \Q$,
  where $p$ and $q$ are relatively prime. By the Euclidean
  algorithm, we can find $x, y \in \Z$ such that $px + qy = 1$. Let
  $\gamma_\alpha = \abcd{p}{-y}{q}{x}$. Note that $\gamma_\alpha \in
  \SL_2(\Z)$ and $\gamma_{\alpha}(\infty) = \alpha$. Also let
  $\gamma_\infty$ be the identity map on $\P^1(\Q)$. Now
  $\gamma_\beta^{-1}$ sends $\beta$ to $\infty$ so we have
  $\gamma_\alpha \circ \gamma_\beta^{-1}$ which sends $\alpha$ to
  $\beta$.

\end{enumerate}

\section{Chapter~\ref{ch:level1}}

\begin{enumerate}
\item \exref{ch:level1}{ex:zetaval}.
We have
$$\zeta(26) = \frac{1315862 \cdot \pi^{26}}{11094481976030578125}.$$
Variation: Compute $\zeta(28)$.

\item \exref{ch:level1}{ex:odd_bernoulli}. Omitted.

\item \exref{ch:level1}{ex:expeis}.
\begin{align*}
 E_8 &=  -\frac{B_8}{16} + q + \sum_{n=2}^{\infty} \sigma_7(n) q^n\\
     &= \frac{1}{480} + q + 129q^{2} + 2188q^{3} + \cdots.
   \end{align*}
Variation: Compute $E_{10}$.


\item \exref{ch:level1}{ex:g4g6quo}. Omitted.

\item  \exref{ch:level1}{ex:vm}.
We have $d=\dim S_{28} = 2$.  A choice of
$a,b$ with $4a+6b\leq 14$ and $4a+6b \con 4\pmod{12}$
is $a=1, b=0$.
A basis for $S_{28}$ is then 
\begin{align*}
\qquad \qquad g_1 &=  \Delta F_6^{2(2-1) + 0} F_4 =q - 792q^{2} - 324q^{3} + 67590208q^{4} + \cdots,\\
g_2 &=  \Delta^2 F_6^{2(2-2)+0} F_4 = q^{2} + 192q^{3} - 8280q^{4} + \cdots.
\end{align*}
The Victor Miller basis is then
\begin{align*}
  f_1 &= g_1 +729g_2 = q + 151740q^{3} + 61032448q^{4} + \cdots, \\
  f_2 &= g_2 = q^{2} + 192q^{3} - 8280q^{4} + \cdots.
\end{align*}
Variation: Compute the Victor Miller basis for $S_{30}$.

\item \exref{ch:level1}{ex:elkies28}.
From the previous exercise we have $f = \Delta^2 F_4$.
Then
\begin{align*}
\qquad \qquad f = \Delta^2 F_4 &= \left(\frac{F_4^3 - F_6^2}{-1728}\right)^2 \cdot F_4 \\
      &= \left(\frac{\left(-\frac{8}{B_4} E_4\right)^3 - 
              \left(-\frac{12}{B_6} E_6\right)^2}{-1728}\right)^2
           \cdot \left(-\frac{8}{B_4} E_4\right) \\
      &= 5186160 E_4E_6^{4} - 564480000 E_4^{4}E_6^{2} + 15360000000 E_4^{7}.
\end{align*}


\item \exref{ch:level1}{ex:fkint}.  
No, it is not always integral.
For example, for $k=12$, the coefficient of $q$
is $-2\cdot 12/B_{12} = 65520/691 \not\in\Z$.
Variation: Find, with proof, the set of all $k$ such that the normalized
series $F_k$ {\em is} integral (use that $B_k$ is eventually very
large compared to $2k$).

\item \exref{ch:level1}{ex:vm32t2}. 
We compute the Victor Miller basis to precision 
great enough to determine $T_2$. This means
we need up to $O(q^5)$. 
\begin{align*}
 f_0 &= 1 + 2611200q^{3} + 19524758400q^{4} + \cdots,\\
 f_1 &= q + 50220q^{3} + 87866368q^{4} +\cdots,\\
 f_2 &= q^{2} + 432q^{3} + 39960q^{4} + \cdots.
\end{align*}
Then the matrix of $T_2$ on this basis is 
$$
\left(\begin{array}{rrr}
2147483649&0&19524758400\\
0&0&2235350016\\
0&1&39960
\end{array}\right).
$$
(The rows of this matrix are the linear combinations that
give the images of the $f_i$ under $T_2$.)
This matrix has characteristic polynomial
$$
   (x - 2147483649) \cdot (x^{2} - 39960x - 2235350016).
$$


    
\end{enumerate}

\section{Chapter~\ref{ch:modform2}}

\begin{enumerate}

\item \exref{ch:modform2}{ex:x0n}. 
Write $g=\abcd{a}{b}{c}{d}$, so 
$\lambda' = \frac{a\lambda + b}{c\lambda + d}$.
Let $f$ be the isomorphism $\C/\Lambda \to \C/\Lambda'$
given by $f(z) = z/(c\lambda + d)$.
We have
$$
 \qquad\quad f\left(\frac{1}{N}\right) = 
    \frac{1}{N(c\lambda + d)} = \frac{a}{N} - \frac{c}{N} \cdot \frac{a \lambda + b}{c\lambda + d}
   \cong \frac{a}{N} \pmod{\Z + \Z \lambda'},
$$
where the second equality can be verified easily by expanding out each side,
and for the congruence we use that $N\mid c$.
Thus the subgroup of $\C/\Lambda$ generated by $\frac{1}{N}$ is
taken isomorphically to the subgroup of $\C/\Lambda'$ generated
by $\frac{1}{N}$.

\item \exref{ch:modform2}{ex:modsymzero}. 
For any integer $r$, 
we have $\abcd{1}{r}{0}{1}\in \Gamma_0(N)$,
so $\{0,\infty\} = \{r, \infty\}$.
Thus
$$
  \qquad  0 = \{0,\infty\} - \{0,\infty\} = 
           \{n,\infty\} - \{m,\infty\} = 
               \{n,\infty\} + \{\infty, m\} = \{n,m\}.
$$

\item \exref{ch:modform2}{ex:bij}.
\begin{enumerate}
\item $(0:1), (1:0), (1:1), \ldots, (1,p-1)$.
\item $p+1$.
\item See \cite[Prop.~2.2.1]{cremona:algs}.
\end{enumerate}

\item \exref{ch:modform2}{ex:intermsofmanin}.
We start with $b=4$, $a=7$.  Then $4\cdot 2 \con 1\pmod{7}$.
Let $\delta_1 = \abcd{4}{1}{7}{2}\in\SL_2(\Z)$.  Since 
$\delta_1\in \Gamma_0(7)$,
we use the right coset representative $\abcd{1}{0}{0}{1}$ and see that
$$
   \{0,4/7\} = \{0,1/2\} + \abcd{1}{0}{0}{1}\{0,\infty\}.
$$
Repeating the process, we have $\delta_2 = \abcd{1}{1}{2}{0}$, which is
in the same coset at $\abcd{0}{-1}{1}{\hfill0}$.  Thus
$$
\{0,1/2\} = \abcd{0}{6}{1}{0} \{0,\infty\} + \{0,0\}.
$$
Putting it together gives 
$$
  \{0,4/7\} = \abcd{1}{0}{0}{1} \{0,\infty\} + \abcd{0}{6}{1}{0}\{0,\infty\}
             = [(0,1)] + [(1,0)].
$$

\item \exref{ch:modform2}{ex:heckes2g03}.
%soln from R. Bradshaw.
\begin{enumerate}
\item Coset representatives for $\Gamma_0(3)$ in $\SL_2(\Z)$ are 
$$
\mtwo{ 1 }{ 0  }{ 0  }{ 1 },\quad
\mtwo{ 1 }{ 0 }{ 1  }{ 1 },\quad
\mtwo{ 1 }{ 0  }{ 2  }{ 1 },\quad
\mtwo{ 0 }{ -1  }{ 1  }{ 0 },
$$
which we refer to below as $[r_0], [r_1], [r_2], $ and $[r_3]$, respectively. 

\item 
In terms of representatives we have 
%
% T := Matrix([[1,-1],[1,0]]);
% S := Matrix([[0,-1],[1,0]]);
% r[0] := Matrix([[1,0],[0,1]]);
% r[1] := Matrix([[1,0],[1,1]]);
% r[2] := Matrix([[1,0],[2,1]]);
% r[3] := Matrix([[0,-1],[1,0]]);

%for i from 0 to 3 do  
%A := multiply(r[i],T^2);
%printf("r[%d] * X  = ", i);
%print(evalm(A));
%for j from 0 to 3 do
%if ((multiply(A, r[j]^(-1))[2,1] mod 3)=0) then printf("in class r[%d]\n", j);
%printf("r[%d] * X * r[%d]^(-1)\n", i, j);
%print(evalm(multiply(A, r[j]^(-1))));
%fi;
%od;
%od:
%
$$
\begin{array}{ccc}
& [r_0]+[r_3]=0, & [r_0]+[r_3]+[r_2]=0, \\
& [r_1]+[r_2]=0, & [r_1]+[r_1]+[r_1]=0, \\
& [r_2]+[r_1]=0, & [r_2]+[r_0]+[r_3]=0, \\
& [r_3]+[r_0]=0, & [r_3]+[r_2]+[r_0]=0.
\end{array}
$$

\item By the first three relations we have $[r_2] = [r_1] = 0 = 0[r_0]$ and $[r_3] = -1[r_0]$. 

\item \begin{align*}
\qquad\quad T_2([r_0]) & =  
[r_0] \abcd{1}{0}{0}{2}
+ [r_0]\abcd{2 }{ 0 }{ 0 }{ 1}
+ [r_0]\abcd{2 }{ 1 }{ 0 }{ 1}
+ [r_0]\abcd{1 }{ 0 }{ 1 }{ 2} \\
& = 
\left[ \abcd{1 }{ 0 }{ 0 }{ 2} \right]
+ \left[ \abcd{2 }{ 0 }{ 0 }{ 1} \right]
+ \left[ \abcd{2 }{ 1 }{ 0 }{ 1} \right]
+ \left[ \abcd{1 }{ 0 }{ 1 }{ 2} \right] \\
& =  [r_0] + [r_0] + [r_0] + [r_2] \\
& =  3[r_0].
\end{align*}

\end{enumerate}


\end{enumerate}

\section{Chapter~\ref{ch:dirichlet}}

\begin{enumerate}
\item \exref{ch:dirichlet}{ex:notdir}.
Suppose~$f$ is a Dirichlet character with modulus~$N$.
Then $-1=f(-1) = f(-1+N)=1$, a contradiction.

\item \exref{ch:dirichlet}{ex:cyclic}.
\begin{enumerate}
\item Any finite subgroup of the multiplicative group of a field
is cyclic (since the number of roots of a polynomial over a field
is at most its degree), so $(\Z/p\Z)^*$ is cyclic.  Let $g$ be an
integer that reduces to a generator of $(\Z/p\Z)^*$.
Let $x=1+p \in (\Z/p^n\Z)^*$; by the binomial theorem 
$$\qquad\qquad x^{p^{n-2}} = 1 + p^{n-2} \cdot p + \cdots
         \con 1 + p^{n-1} \not\con 0 \pmod{p^{n}},
$$
so $x$ has order $p^{n-1}$. Since $p$ is odd, $\gcd(p^{n-1},p-1) = 1$,
so $xg$ has order $p^{n-1}\cdot (p-1) = \vphi(p^n)$; hence
$(\Z/p^n\Z)^*$ is cyclic.  
\item By the binomial theorem $(1+2^2)^{2^{n-3}} \not\con 1 \pmod{2^n}$,
so $5$ has order $2^{n-2}$ in $(\Z/2^n\Z)^*$, and clearly $-1$ has
order $2$. Since $5\con 1\pmod{4}$, $-1$ is not a power of $5$ in 
$(\Z/2^n\Z)^*$.  Thus the subgroups $\langle -1 \rangle$
and $\langle 5 \rangle$ have trivial intersection.  
The product of their orders is $2^{n-1} = \vphi(2^n) = \#(\Z/2^n\Z)^*$,
so the claim follows. 
\end{enumerate}

\item\exref{ch:dirichlet}{ex:orderalg}.
Write $n=\prod p_i^{e_i}$.
The order of $g$ divides $n$, so 
the condition implies that $p_i^{e_i}$ divides the order of $g$
for each $i$.   Thus the order of $g$ is divisible by the
least common multiple of the $p_i^{e_i}$, i.e., by $n$.


\item\exref{ch:dirichlet}{ex:dlogadd}.
\begin{enumerate}
\item 
The bijection given by
$1+p^{n-1}a\pmod{p^n} \mapsto a\pmod{p}$
 is a homomorphism
since $$\hspace{6em}(1+p^{n-1}a)(1+p^{n-1}b) \con 1 + p^{n-1}(a+b) \pmod{p^n}.$$
\item We have an exact sequence
$$
 \qquad\qquad  1 \to 1 + p \Z/p^{n}\Z \to (\Z/p^n\Z)^* \to (\Z/p\Z)^*\to 1,
$$
so it suffices to solve the discrete log problem in the kernel
and cokernel. We prove by induction on $n$ that we can solve
the discrete log problem
in the kernel easily (compared to known methods for solving
the discrete log problem in $(\Z/p\Z)^*$).
We have an exact sequence
$$
\qquad\qquad  1 \to 1 + p^{n-1} \Z/p^{n}\Z \to (\Z/p^n\Z)^*  \to 
             (\Z/p^{n-1}\Z)^* \to 1.$$
The
first part of this problem shows that we can solve the
discrete log problem in the kernel, and by induction
we can solve it in the cokernel.  This completes
the proof.
\end{enumerate}

\item \exref{ch:dirichlet}{ex:cond2}.
If $\eps(5)=1$, then since $\eps$ is 
nontrivial, \exref{ch:dirichlet}{ex:cyclic} implies
that $\eps$ factors through $(\Z/4\Z)^*$, hence has
conductor $4 = 2^{1 + 1}$, as claimed. 
If $\eps(5)\neq 1$, then again from \exref{ch:dirichlet}{ex:cyclic}
we see that if $\eps$ has order~$r$, then~$\eps$ 
factors through $(\Z/2^{r+2}\Z)^*$ but nothing smaller. 

\item \exref{ch:dirichlet}{ex:charmod5}.
\begin{enumerate}
\item Take $f=x^2+2$.
\item The element $2$ has order $4$.
\item A minimal generator for $(\Z/25\Z)^*$ is $2$,
and the characters are $[1]$, $[2]$, $[3]$, $[4]$.
\item Each of the four Galois orbits has size $1$.
\end{enumerate}


\end{enumerate}

\section{Chapter~\ref{ch:eisen}}

\begin{enumerate}
\item \exref{ch:eisen}{ex:diag}.
The eigenspace $E_{\lambda}$ of $A$ with eigenvalue $\lambda$ is
preserved by~$B$, since if $v\in E_{\lambda}$, then 
$$ABv = BAv = B(\lambda v) = \lambda B v.$$
Because $B$ is diagonalizable, its minimal polynomial
equals its characteristic polynomial; hence the
same is true for the restriction of $B$ to $E_{\lambda}$,
i.e., the restriction of $B$ is diagonalizable.  
Choose basis for all $E_{\lambda}$ so that the
restrictions of $B$ to these eigenspaces
is diagonal with respect to these bases. Then
the concatenation of these bases is a basis that
simultaneously diagonalizes $A$ and $B$.

\item \exref{ch:eisen}{ex:bern_triv}.
When $\eps$ is the trivial character, the
$B_{k,\eps}$ are defined by
$$
  \qquad \sum_{a=1}^{1} \frac{\eps(a) x e^{ax}}{e^{x}-1}
   = \frac{x e^x}{e^x - 1} = x + \frac{x}{e^x - 1}
   = \sum_{k=0}^{\infty} B_{k,\eps} \frac{x^k}{k!}.
$$
Thus $B_{1,\eps} = 1 + B_1 = \frac{1}{2}$,
and for $k>1$, we have $B_{k,\eps} = B_k$.

\item \exref{ch:eisen}{ex:gbparity}. Omitted.

\item \exref{ch:eisen}{ex:dime3}.
The  Eisenstein series in our
basis for $E_3(\Gamma_1(13))$ are of 
the form $E_{3,1,\eps}$ or $E_{3,\eps,1}$
with $\eps(-1)=(-1)^3=-1$.
There are six characters~$\eps$
with modulus $13$ such that  $\eps(-1)=-1$,
and we have the two series $E_{3,1,\eps}$ and
$E_{3,\eps,1}$ associated to each of these.
This gives a dimension of $12$.


\end{enumerate}

\section{Chapter~\ref{ch:dim}}
\begin{enumerate}
\item \exref{ch:dim}{ex:mu0n}.
\begin{enumerate}
\item By Proposition~\ref{prop:p1bij}, we have 
$[\SL_2(\Z):\Gamma_0(N)] = \#\P^1(\Z/N\Z)$.  
By the Chinese Remainder Theorem, 
$$\#\P^1(\Z/N\Z) = \prod_{p\mid N} \#\P^1(\Z/p^{\ord_p(N)}\Z).$$
So we are reduced to computing $\#\P^1(\Z/p^{\ord_p(N)}\Z)$.
We have $(a,b)\in (\Z/p^n\Z)^2$ with $\gcd(a,b,p)=1$
if and only if $(a,b)\not\in (p\Z/p^n\Z)^2$, so
there are $p^{2n} - p^{2(n-1)}$ such pairs. 
The unit group $(\Z/p^n\Z)^*$ has order $\vphi(p^n) = p^n - p^{n-1}$.
It follows that 
$$
  \#\P^1(\Z/N\Z) = \frac{p^{2n} - p^{2(n-1)}}{p^n - p^{n-1}}
   = p^{n} + p^{n-1}.
$$
\item Omitted.
\end{enumerate}

\item \exref{ch:dim}{ex:formsl2z}.  Omitted.

\item \exref{ch:dim}{ex:newmatch}. Omitted.

\item \exref{ch:dim}{ex:suma4}. Omitted.

\item \exref{ch:dim}{ex:sagedim}. See the source
code to \sage.  

\end{enumerate}


\section{Chapter~\ref{ch:linalg}}
\begin{enumerate}
\item \exref{ch:linalg}{ex:matrixwithkernel}.
Take a basis of $W$ and let $G$ be the matrix
 whose rows are these basis elements. Let
 $B$ be the row echelon form of $G$. 
After a permutation $p$ of columns, we may
 write $B = p_i(I | C)$, where $I$ is the identity matrix.
 The matrix $A = p^{-1}(-C^t | I)$, where $I$ is a different
 sized identity matrix, has the property that
 $W = \Ker(A)$.

\item \exref{ch:linalg}{ex:rrefmod}.
The answer is no.  For example if $A=nI$ is $n$ times the
identity matrix and if $p \mid n$, then $\rref(A\pmod{p})=0$
but $\rref(A)\pmod{p}=I$.

\item \exref{ch:linalg}{ex:rrefsmodp}.
Let $T=\prod E_i$ be an invertible matrix such that $TA = E$ is
in (reduced) echelon form and the $E_i$ are elementary matrices,
i.e., the result of applying an elementary row operation to the
identity matrix.  If $p$ is a prime
that does not divide any of the nonzero numerators
or denominators of the entries of $A$ and any $E_i$, then 
$\rref(A \pmod{p}) = \rref(A) \pmod{p}$.  This is
because $E\pmod{p}$ is in echelon form and $A\pmod{p}$
can be transformed to $E\pmod{p}$ via a series of
elementary row operations modulo~$p$.

\item \exref{ch:linalg}{ex:linalg1to9}.
\begin{enumerate}
\item The echelon form (over $\Q$) is 
$$\left(\begin{array}{rrr}
1&0&-1\\
0&1&2\\
0&0&0
\end{array}\right).$$
\item The kernel is the $1$-dimensional span of 
$\left(1,-2,1\right)$.
\item The characteristic polynomial is $x \cdot (x^{2} - 15x - 18)$.
\end{enumerate}

%\item \exref{ch:linalg}{ex:transmat}.

\item \exref{ch:linalg}{ex:hnf}.
\begin{enumerate}
\item The answer is given in the problem.
\item See \cite[\S2.4]{cohen:course_ant}. 
\end{enumerate}

\end{enumerate}

\section{Chapter~\ref{ch:modsym}}
\begin{enumerate}
\item \exref{ch:modsym}{ex:surjredmod}.  
Using the Chinese Remainder Theorem we immediately reduce
to proving the statement when both $M=p^r$ and $N=p^s$ are powers
of a prime $p$. Then $[a]\in(\Z/p^s\Z)^*$ is represented by an
integer $a$ with $\gcd(a,p)=1$.  That same integer $a$
defines an element of $(\Z/p^r\Z)^*$ that reduces modulo
$p^s$ to $[a]$.



\item \exref{ch:modsym}{ex:surjred}. See \cite[Lemma 1.38]{shimura:intro}.


\item \exref{ch:modsym}{ex:m3g13}.
Coset representatives for $\Gamma_1(3)$ are in bijection with $(c,d)$ where 
$c,d \in \Z/3\Z$ and gcd$(c,d,N)=1$, so the following are representatives:
$$ 
\abcd{1 }{0}{0}{1},\,
\abcd{2 }{0}{0}{2}, \,
\abcd{0 }{2}{1}{0},\,
\abcd{1 }{0}{1}{1},\,
\abcd{2 }{0}{1}{2},\,
\abcd{1 }{1}{2}{0},\,
\abcd{1 }{0}{2}{1},\,
\abcd{2 }{0}{2}{2},
$$
which we call $r_1, \ldots, r_8$, respectively.
Now our Manin symbols are of the form $[X,r_i]$ 
and $[Y, r_i]$ for $1 \le i \le 8$ 
modulo the relations 
$$x+x\sigma = 0, \quad x+x\tau+x\tau^2 = 0, \quad\text{ and } x-xJ = 0.$$
First, note that $J$ acts trivially on Manin symbols of odd weight because 
it sends $X$ to $-X$, $Y$ to $-Y$ and $r_i$ to $-r_i$, so $$[z,g]J = [-z,-g] = [z,g].$$ Thus the last relation is trivially true. 

Now $\sigma^{-1}X = -Y$ and $\sigma^{-1}Y = X$. 
Also $\tau^{-1}X = -Y, \tau^{-1}Y = X-Y, \tau^{-2}X = -X+Y$ 
and $\tau^{-2}Y = -X$. 

The first relation on the first symbol says that
$$[X, r_1] = -[-Y, r_3] = [Y, r_3]$$
and the second relation tells us that
$$[X, r_1] + [-Y, r_5] + [-X+Y, r_6] = 0.$$



\item \exref{ch:modsym}{ex:pairingwd}.
  Let $f \in S_k(\Gamma)$ and $g \in \Gamma$. All that remains to be
  shown is that this pairing respects the relation $x=xg$ for all
  modular symbols $x$.  By linearity it suffices to show the invariance of
  $\langle f, X^{k-i-2}Y^i\{\alpha, \beta\} \rangle$.  We have
\begin{align*}
\qquad&\left\langle f, (X^{k-2-i}Y^i \{\alpha, \beta\})g^{-1} \right\rangle\\
& =  \left \langle f, (aX+bY)^{k-i-2}(cX+dY)^i \{g^{-1}(\alpha), g^{-1}(\beta)\} \right \rangle \\
& =  \int_{g^{-1}(\alpha)}^{g^{-1}(\beta)} f(z) (az+b)^{k-i-2}(cz+d)^i \; dz \\
& =  \int_{g^{-1}(\alpha)}^{g^{-1}(\beta)} f(z) \frac{(az+b)^{k-i-2}}{(cz+d)^{k-i-2}}(cz+d)^{k-2} \; dz \\
& =  \int_{g^{-1}(\alpha)}^{g^{-1}(\beta)} f(z) \, g(z)^{k-i-2} (cz+d)^{k-2} \; dz \\
& =  \int_{\alpha}^{\beta} f(g^{-1}(z)) \, g(g^{-1}(z))^{k-i-2} (cg^{-1}(z)+d)^{k-2} \; d(g^{-1}(z)) \\
& =  \int_{\alpha}^{\beta} f(g^{-1}(z)) \, z^{k-i-2} (cg^{-1}(z)+d)^{k-2} \; (cg^{-1}(z)+d)^2 \; dz \\
& =  \int_{\alpha}^{\beta} f(z) \, z^{k-i-2} \; dz \\
& =  \left \langle f, X^{k-i-2}Y^i \{ \alpha , \beta \} \right \rangle,
\end{align*}
where the second to last simplification is due to invariance under
$[g]_k$, i.e.,
$$\qquad f(g^{-1}(z)) = f^{[g]_k}(g^{-1}(z)) = (cg^{-1}(z)+d)^{-k} f(g(g^{-1}(z))).$$
(The proof for $f \in \overline{S}_k(\Gamma)$ works in
 exactly the same way.) 

\item \exref{ch:modsym}{ex:etaconj}.
\begin{enumerate}
\item Let $\eta = {\scriptsize \left(\begin{array}{rr}-1 & 0 \\ 0 & 1
      \end{array}\right)} $. For any $\gamma = {\scriptsize
    \left(\begin{array}{cc}a & b \\c & d \end{array}\right)}$ we have
$$\qquad\qquad \gamma \eta =  {\scriptsize  \left(\begin{array}{rr}-a & b \\ -c & d \end{array}\right)},~~~~
\eta \gamma = {\scriptsize \left(\begin{array}{rr}-a & -b \\ c & d
    \end{array}\right)},~~~~ \mbox{and}~~~~~ \eta \gamma \eta =
{\scriptsize \left(\begin{array}{rr} a & -b \\ -c & d
    \end{array}\right)}.
 $$
 First, if $\gamma \in SL_2(\Z)$, then $\eta \gamma \eta \in GL_2(\Z)$
 and $$\det (\eta \gamma \eta) = \det \eta \det \gamma \det \eta =
 (-1)(1)(-1) = 1$$ so $\eta \gamma \eta \in SL_2(\Z)$. As $\eta^2 = 1$,
 conjugation by $\eta$ is self-inverse, so it must be a bijection.

 Now if $\gamma \in \Gamma_0(N)$, then $c \equiv 0$ (mod $N$), so $-c
 \equiv 0$ (mod $N$), and so $\eta \gamma \eta \in \Gamma_0(N)$. Thus
 $\eta\Gamma_0(N)\eta = \Gamma_0(N)$.

 If $\gamma \in \Gamma_1(N)$, then $-c \equiv 0$ (mod $N$) as before
 and also $a \equiv d \equiv 1$ (mod $N$), so $\eta \gamma \eta \in
 \Gamma_1(N)$. Thus $\eta\Gamma_1(N)\eta = \Gamma_1(N)$.

\item Omitted.

\end{enumerate}


\end{enumerate}


\section{Chapter~\ref{ch:modformcomp}}
\begin{enumerate}
\item \exref{ch:modformcomp}{ex:g0g1quo}.
Consider the surjective homomorphism $$r:\SL_2(\Z)\to \SL_2(\Z/N\Z).$$
Notice that $\Gamma_1(N)$ is the exact inverse image of the subgroup $H$ of
matrices of the form $\abcd{1}{*}{0}{1}$ and $\Gamma_0(N)$
is the inverse image of the subgroup~$T$ of upper triangular matrices.
It thus suffices to observe that $H$ is normal in $T$, which is clear.
Finally, the quotient $T/H$ is isomorphic to the group of diagonal
matrices in $\SL_2(\Z/N\Z)^*$, which is isomorphic to $(\Z/N\Z)^*$.

\item \exref{ch:modformcomp}{ex:dbdinhecke}.
It is enough to show $\dbd{p}\in\Z[\ldots,T_n,\ldots]$ for
primes~$p$, since each $\dbd{d}$ can be written in terms of the
$\dbd{p}$. Since $p\nmid N$, we have that
\[
  T_{p^2} = T_p^2 - \dbd{p}p^{k-1},
\]
so $$\dbd{p}p^{k-1} = T_p^2 - T_{p^2}.$$
By Dirichlet's theorem on primes in arithmetic progression, 
there is a prime~$q\neq p$ congruent
to~$p$ mod $N$. Since $p^{k-1}$ and $q^{k-1}$ are relatively
prime, there exist integers~$a$ and~$b$ such that $a p^{k-1} + b
q^{k-1} = 1$. Then
$$
\qquad\dbd{p}=\dbd{p}(a p^{k-1} + b q^{k-1})
       = a({T_p}^2-T_{p^2}) + b({T_q}^2-T_{q^2})
       \in \Z[\ldots, T_n,\ldots].
$$

\item \exref{ch:modformcomp}{exer:onlyold}.
Take $N=33$. The space $S_2(\Gamma_0(33))$ is a direct
sum of the two old subspaces coming from $S_2(\Gamma_0(11))$
and the new subspace, which has dimension $1$.
If $f$ is a basis for $S_2(\Gamma_0(11))$ and
$g$ is a basis for $S_2(\Gamma_0(33))_{\new}$,
then $\alpha_1(f), \alpha_3(f), g$ is a basis
for $S_2(\Gamma_0(33))$ on which all Hecke
operators $T_n$, with $\gcd(n,33)=1$, have
diagonal matrix.
However, the operator
$T_3$ on $S_2(\Gamma_0(33))$ does
not act as a scalar on $\alpha_1(f)$,
so it cannot be in the ring generated
by all operators $T_n$ with $\gcd(n,33)=1$.

\item \exref{ch:modformcomp}{ex:31}. Omitted.

\end{enumerate}

\section{Chapter~\ref{ch:periods}}
\begin{enumerate}
\item \exref{ch:periods}{ex:intexp}. Hint: Use either repeated
integration by parts or a change of variables that relates the integral
to the $\Gamma$ function.

\item \exref{ch:periods}{ex:funceqn}. See \cite[\S2.8]{cremona:algs}.


\item \exref{ch:periods}{ex:qvs}.
\begin{enumerate}
\item Let $f = \sqrt{-1} \sum q^n$.  Then $d=2$, but the nontrivial 
conjugate of $f$ is $-f$, so $V_f$ has dimension $1$.
\item Choose $\alpha \in K$ such that $K=\Q(\alpha)$.
 Write 
\begin{equation}\label{eqn:fidag}
\qquad\qquad f=\sum_{i=0}^{d-1} \alpha^i g_i
\end{equation}
with $g_i \in \Q[[q]]$.
Let $W_g$ be the $\Q$-span of the $g_i$, and let
$W_f = V_f \cap \Q[[q]]$. 
 By considering the $\Gal(\Qbar/\Q)$ conjugates of (\ref{eqn:fidag}),
we see that the Galois conjugates of $f$ are in the $\C$-span 
of the $g_i$, so
\begin{equation}\label{eqn:dcvfqw}
\qquad\qquad d = \dim_\C V_f \leq \dim_\Q W_g.
\end{equation}
Likewise, taking the above modulo $O(q^n)$ for any~$n$, 
we obtain a matrix equation 
$$
\qquad\qquad  F = A G,
$$
where the columns of $F$ are the $\Gal(\Qbar/\Q)$-conjugates of
$f$, the matrix $A$ is the Vandermonde matrix corresponding to
$\alpha$ (and its $\Gal(\Qbar/\Q)$ conjugates),  and~$G$ has
columns~$g_i$.  Since~$A$ is a Vandermonde matrix, it is 
invertible, so $A^{-1}F = G$.  Taking the limit as $n$ goes
to infinity, we see that each $g_i$ is a linear combination
of the $f_i$, hence an element of $V_f$.   
Thus $W_g \subset W_f$, so (\ref{eqn:dcvfqw})
implies that $\dim_\Q W_f \geq d$.  But $W_f \tensor_\Q \C \subset V_f$
so finally
$$
\qquad\qquad  d\leq \dim_\Q W_f = \dim_\C (W_f \tensor_\Q \C) \leq \dim_\C V_f = d.
$$

\end{enumerate}

\item \exref{ch:periods}{ex:11}. See the appendix to Chapter II in
  \cite{cremona:algs}, where this example is worked out in complete
  detail.  

\end{enumerate}

%%% Local Variables: 
%%% mode: latex
%%% TeX-master: t
%%% End: 
